% ============================================================
% TEMPLATE BÁO CÁO NGHIÊN CỨU KHOA HỌC KỸ THUẬT
% Dùng làm khung hướng dẫn cho nhiều đề tài khác nhau.
% Biên dịch bằng XeLaTeX hoặc LuaLaTeX (khuyến nghị XeLaTeX).
%
% CÁCH DÙNG:
% 1. Thay nội dung trong [ngoặc vuông] bằng thông tin dự án.
% 2. Các đoạn bắt đầu bằng "HƯỚNG DẪN:" là ghi chú cho người viết;
%    xóa chúng trước khi nộp báo cáo chính thức.
% 3. Xóa các mục không phù hợp và thêm mục cần thiết cho đề tài.
% 4. Không bịa kết quả, số liệu, tài liệu tham khảo hoặc kết luận.
% ============================================================

\documentclass[12pt,a4paper]{report}

% ---------- Ngôn ngữ và phông chữ (XeLaTeX/LuaLaTeX) ----------
\usepackage{fontspec}
\usepackage{polyglossia}
\setdefaultlanguage{vietnamese}
\setmainfont{Times New Roman} % Nếu không có, thử TeX Gyre Termes
\setsansfont{Arial}           % Có thể đổi thành TeX Gyre Heros
\setmonofont{Courier New}

% ---------- Trình bày trang ----------
\usepackage[a4paper,top=2.5cm,bottom=2.5cm,left=3cm,right=2cm]{geometry}
\usepackage{setspace}
\onehalfspacing
\usepackage{indentfirst}
\setlength{\parindent}{1.25cm}
\setlength{\parskip}{0pt}
\usepackage{microtype}

% ---------- Hình, bảng, công thức ----------
\usepackage{graphicx}
\usepackage{float}
\usepackage{booktabs}
\usepackage{tabularx}
\usepackage{array}
\usepackage{amsmath,amssymb}
\usepackage{caption}
\usepackage{subcaption}
\captionsetup{font=small,labelfont=bf}
\usepackage{siunitx}

% ---------- Mục lục, liên kết và trích dẫn ----------
\usepackage{hyperref}
\hypersetup{
  colorlinks=true,
  linkcolor=black,
  citecolor=black,
  urlcolor=blue,
  pdftitle={[Tên đề tài]},
  pdfauthor={[Tên tác giả/nhóm tác giả]}
}
\usepackage[numbers,sort&compress]{natbib}
\bibliographystyle{unsrtnat}

% ---------- Định dạng tiêu đề ----------
\usepackage{titlesec}
\titleformat{\chapter}[hang]
  {\normalfont\bfseries\Large}{\thechapter.}{0.7em}{}
\titlespacing*{\chapter}{0pt}{-10pt}{18pt}
\titleformat{\section}
  {\normalfont\bfseries\large}{\thesection}{0.7em}{}
\titleformat{\subsection}
  {\normalfont\bfseries\normalsize}{\thesubsection}{0.7em}{}

% ---------- Lệnh hỗ trợ người viết ----------
\newcommand{\huongdan}[1]{%
  \par\smallskip\noindent
  \textit{\textbf{HƯỚNG DẪN NGƯỜI VIẾT:} #1}%
  \par\smallskip}
\newcommand{\dien}[1]{\textbf{[#1]}}
\newcommand{\tukhoa}[1]{\textbf{#1}}

\begin{document}

% ============================================================
% BÌA
% ============================================================
\begin{titlepage}
  \centering
  {\large \textbf{[TÊN CƠ QUAN / TRƯỜNG / ĐƠN VỊ]}\par}
  \vspace{1.5cm}
  {\large \textbf{[TÊN CUỘC THI / CHƯƠNG TRÌNH NGHIÊN CỨU]}\par}
  \vspace{2.2cm}
  {\Large\bfseries BÁO CÁO NGHIÊN CỨU KHOA HỌC KỸ THUẬT\par}
  \vspace{1cm}
  {\huge\bfseries [TÊN ĐỀ TÀI]\par}
  \vspace{1.5cm}
  \begin{flushleft}
    \textbf{Lĩnh vực:} [Lĩnh vực dự thi/nghiên cứu]\\[0.4cm]
    \textbf{Tác giả:} [Họ tên tác giả hoặc các thành viên]\\[0.4cm]
    \textbf{Người hướng dẫn:} [Nếu có]\\[0.4cm]
    \textbf{Đơn vị:} [Trường/lớp/đơn vị]
  \end{flushleft}
  \vfill
  {\large [Địa điểm] -- [Năm]\par}
\end{titlepage}

% ============================================================
% MỤC LỤC
% ============================================================
\pagenumbering{roman}
\tableofcontents
\clearpage
\pagenumbering{arabic}

% ============================================================
\chapter{Giới thiệu vấn đề nghiên cứu}
% Mục tiêu chương: giúp người đọc hiểu vấn đề, tầm quan trọng,
% nghiên cứu trước đây và khoảng trống mà dự án muốn giải quyết.
% ============================================================

\section{Dẫn nhập}

\huongdan{Mở đầu từ bối cảnh rộng của lĩnh vực, sau đó thu hẹp dần
vào một vấn đề cụ thể. Giải thích vấn đề ảnh hưởng đến ai hoặc điều gì,
vì sao đáng nghiên cứu, và các phương pháp hiện tại còn hạn chế ở đâu.
Dùng dẫn chứng đáng tin cậy và trích dẫn nguồn. Tránh chỉ khẳng định
rằng đề tài ``rất quan trọng'' mà không có bằng chứng.}

[Viết phần bối cảnh và giới thiệu lĩnh vực.]

\subsection{Vấn đề cần giải quyết}

\huongdan{Mô tả cụ thể hiện tượng, khó khăn hoặc nhu cầu chưa được đáp ứng.
Nếu có thể, đưa ví dụ minh họa, dữ liệu nền, tình huống thực tế hoặc hình
ảnh thể hiện hạn chế của phương pháp hiện tại.}

[Trình bày vấn đề cụ thể và ví dụ minh họa.]

\subsection{Mục tiêu nghiên cứu}

\huongdan{Nêu mục tiêu chính có thể kiểm chứng. Có thể tách thành mục tiêu
tổng quát và các mục tiêu cụ thể. Mục tiêu nên mô tả điều dự án sẽ xây dựng,
đề xuất, đo lường hoặc so sánh.}

Mục tiêu chính của đề tài là [mục tiêu chính].

Các mục tiêu cụ thể:
\begin{itemize}
  \item [Mục tiêu cụ thể 1.]
  \item [Mục tiêu cụ thể 2.]
  \item [Mục tiêu cụ thể 3.]
\end{itemize}

\subsection{Câu hỏi nghiên cứu và/hoặc giả thuyết}

\huongdan{Viết câu hỏi có thể trả lời bằng phương pháp và dữ liệu của dự án.
Nếu lĩnh vực phù hợp, nêu giả thuyết cùng cơ sở dự đoán. Không bắt buộc
phải có giả thuyết nếu nghiên cứu mang tính khám phá hoặc xây dựng hệ thống.}

\begin{enumerate}
  \item [Câu hỏi nghiên cứu 1.]
  \item [Câu hỏi nghiên cứu 2.]
  \item [Câu hỏi nghiên cứu 3, nếu cần.]
\end{enumerate}

\subsection{Phạm vi nghiên cứu}

\huongdan{Nêu rõ đối tượng, dữ liệu, điều kiện thử nghiệm, giới hạn thời gian
và những vấn đề dự án không giải quyết. Phạm vi hợp lý giúp tránh tuyên bố
quá rộng so với bằng chứng thu được.}

[Đối tượng, phạm vi, điều kiện và giới hạn của nghiên cứu.]

\section{Các nghiên cứu có liên quan}

\huongdan{Tổng hợp những công trình, sản phẩm hoặc phương pháp gần với đề tài.
Với mỗi hướng nghiên cứu, trình bày: ý tưởng chính, điểm mạnh, hạn chế và
mối liên hệ với vấn đề của dự án. Không chỉ liệt kê từng bài báo; cần so sánh
để dẫn tới khoảng trống nghiên cứu. Trích dẫn đầy đủ các nguồn đã sử dụng.}

\subsection{Hướng nghiên cứu/phương pháp thứ nhất}

[Giải thích phương pháp, kết quả chính, ưu điểm và hạn chế.] \citep{ref1}

\subsection{Hướng nghiên cứu/phương pháp thứ hai}

[Giải thích phương pháp, kết quả chính, ưu điểm và hạn chế.] \citep{ref2}

\subsection{So sánh và xác định khoảng trống nghiên cứu}

\huongdan{Tóm tắt điều đã được giải quyết và điều còn chưa được giải quyết.
Nêu rõ dự án kế thừa gì, khác biệt ở đâu và đóng góp dự kiến là gì. Tránh
tuyên bố ``chưa ai từng nghiên cứu'' nếu chưa có khảo sát tài liệu đủ rộng.}

[Khoảng trống nghiên cứu và lý do dự án đề xuất cách tiếp cận hiện tại.]

% ============================================================
\chapter{Phương pháp và quá trình thực hiện}
% Có thể đổi tên chương thành "Tóm tắt nghiên cứu", "Thiết kế hệ thống"
% hoặc "Phương pháp nghiên cứu" cho phù hợp lĩnh vực.
% ============================================================

\section{Tổng quan giải pháp đề xuất}

\huongdan{Tóm tắt giải pháp trong một vài đoạn: đầu vào là gì, quy trình chính
gồm những bước nào, đầu ra là gì và giải pháp xử lý hạn chế đã nêu ra ra sao.
Nếu hữu ích, đặt sơ đồ tổng quan ở đây.}

[Giới thiệu ngắn gọn giải pháp đề xuất.]

% Ví dụ chèn hình: thay "figures/overview.png" bằng đường dẫn hình của bạn.
% Nếu chưa có hình, hãy xóa khối figure này.
\begin{figure}[H]
  \centering
  % \includegraphics[width=0.85\textwidth]{figures/overview.png}
  \fbox{\parbox[c][3cm][c]{0.8\textwidth}{\centering
  [Chèn sơ đồ tổng quan giải pháp tại đây]}}
  \caption{Sơ đồ tổng quan của giải pháp đề xuất.}
  \label{fig:overview}
\end{figure}

\section{Cơ sở lý thuyết và công cụ}

\huongdan{Chỉ trình bày kiến thức trực tiếp cần để hiểu giải pháp: thuật toán,
mô hình, định nghĩa, công thức, thư viện, thiết bị hoặc công cụ liên quan.
Phân biệt kiến thức kế thừa từ tài liệu với phần do nhóm tự đề xuất.}

[Trình bày nền tảng lý thuyết, công cụ và lý do lựa chọn.]

\section{Thiết kế phương pháp hoặc kiến trúc hệ thống}

\huongdan{Mô tả đủ chi tiết để người khác hiểu cách giải pháp hoạt động.
Tùy lĩnh vực, chia thành các module, giai đoạn, thành phần, quy trình hoặc
thí nghiệm. Với mỗi thành phần, nêu đầu vào, xử lý và đầu ra.}

\subsection{Thành phần/giai đoạn thứ nhất}

[Chức năng, cách hoạt động, dữ liệu đầu vào và đầu ra.]

\subsection{Thành phần/giai đoạn thứ hai}

[Chức năng, cách hoạt động, dữ liệu đầu vào và đầu ra.]

\subsection{Thuật toán hoặc quy trình thực hiện}

\huongdan{Có thể dùng lưu đồ, giả mã, công thức hoặc các bước đánh số.
Giải thích các ký hiệu trong công thức và các tham số quan trọng. Không
đưa công thức vào chỉ để làm báo cáo trông có vẻ kỹ thuật.}

\begin{enumerate}
  \item [Bước 1: mô tả thao tác.]
  \item [Bước 2: mô tả thao tác.]
  \item [Bước 3: mô tả thao tác.]
\end{enumerate}

% Ví dụ công thức (xóa nếu không phù hợp):
\[
  [\text{Công thức hoặc mô hình toán học của đề tài}]
\]

\section{Dữ liệu, vật liệu hoặc môi trường thực hiện}

\huongdan{Mô tả nguồn dữ liệu/vật liệu, cách thu thập, tiêu chí chọn mẫu,
cách tiền xử lý, thiết bị, phần cứng/phần mềm và môi trường thử nghiệm.
Nêu giấy phép hoặc vấn đề đạo đức/quyền riêng tư nếu có liên quan.}

[Thông tin dữ liệu, vật liệu, thiết bị và môi trường.]

\section{Triển khai sản phẩm hoặc nguyên mẫu}

\huongdan{Nếu dự án tạo ra phần mềm, thiết bị hoặc ứng dụng, mô tả cách
hiện thực hóa giải pháp, kiến trúc triển khai và chức năng chính. Nếu dự án
không tạo sản phẩm, có thể đổi mục này thành quy trình thực nghiệm hoặc
xóa bỏ. Phân biệt phần đã hoàn thành với phần mới ở mức thiết kế.}

[Công nghệ sử dụng, cách triển khai, chức năng và trạng thái hoàn thành.]

% ============================================================
\chapter{Kết quả và đánh giá}
% Chỉ báo cáo kết quả đã thực sự thu được. Không để số liệu ví dụ
% trong bản nộp chính thức.
% ============================================================

\section{Thiết kế đánh giá/thực nghiệm}

\huongdan{Nêu cách kiểm tra giải pháp có đạt mục tiêu hay không. Xác định
đối chứng hoặc phương pháp cơ sở, bộ dữ liệu/mẫu thử, số lần chạy, chỉ số
đánh giá và điều kiện để so sánh công bằng. Nếu có khảo sát người dùng,
mô tả người tham gia, câu hỏi, quy trình và cách tổng hợp phản hồi.}

[Thiết kế thí nghiệm, tiêu chí đánh giá và phương pháp đối chứng.]

\section{Kết quả chính}

\huongdan{Trình bày kết quả bằng bảng, biểu đồ, hình ảnh hoặc ví dụ đại diện.
Mỗi hình/bảng cần có số thứ tự, chú thích và được nhắc tới trong phần văn bản.
Nêu kết quả quan sát được, không chỉ nói rằng phương pháp ``tốt hơn''.}

\begin{table}[H]
  \centering
  \caption{Mẫu bảng tổng hợp kết quả thực nghiệm.}
  \label{tab:results}
  \begin{tabularx}{\textwidth}{Xcc}
    \toprule
    \textbf{Phương pháp/điều kiện} & \textbf{Chỉ số 1} & \textbf{Chỉ số 2} \\
    \midrule
    [Phương pháp cơ sở] & [Kết quả] & [Kết quả] \\
    [Phương pháp đề xuất] & [Kết quả] & [Kết quả] \\
    \bottomrule
  \end{tabularx}
\end{table}

\begin{figure}[H]
  \centering
  % \includegraphics[width=0.8\textwidth]{figures/result.png}
  \fbox{\parbox[c][3cm][c]{0.8\textwidth}{\centering
  [Chèn biểu đồ, ảnh kết quả hoặc minh họa tại đây]}}
  \caption{[Mô tả chính xác nội dung hình và điều kiện tạo ra kết quả.]}
  \label{fig:result}
\end{figure}

\section{Đánh giá định lượng}

\huongdan{Báo cáo chỉ số định lượng phù hợp với bài toán, định nghĩa cách tính
và giải thích ý nghĩa. So sánh với baseline hoặc các điều kiện khác nếu có.
Nếu kết quả biến thiên, trình bày độ lệch, khoảng tin cậy hoặc số lần lặp khi
phù hợp. Không chọn chỉ số chỉ vì nó làm kết quả trông tốt hơn.}

[Chỉ số, kết quả, so sánh và diễn giải.]

\section{Đánh giá định tính hoặc đánh giá bởi con người}

\huongdan{Dùng khi phù hợp với đề tài. Nêu tiêu chí đánh giá, cách chọn người
đánh giá, số lượng mẫu và cách giảm thiên lệch. Có thể đưa ví dụ thành công
và thất bại, phản hồi người dùng hoặc phân tích trường hợp điển hình. Nếu không
thực hiện loại đánh giá này, hãy xóa mục hoặc giải thích vì sao không áp dụng.}

[Kết quả đánh giá định tính/khảo sát và các ví dụ minh họa.]

\section{Thảo luận kết quả}

\huongdan{Giải thích kết quả cho thấy điều gì đối với câu hỏi nghiên cứu.
Nêu kết quả phù hợp hoặc không phù hợp với giả thuyết, các đánh đổi, trường hợp
thất bại và những yếu tố có thể ảnh hưởng. Phân biệt rõ quan sát thực tế với
suy luận của tác giả.}

[Diễn giải ý nghĩa, nguyên nhân, đánh đổi và giới hạn của kết quả.]

\section{Hạn chế của nghiên cứu}

\huongdan{Nêu giới hạn về dữ liệu, cỡ mẫu, tài nguyên, độ tổng quát, thời gian,
độ tin cậy hoặc điều kiện thử nghiệm. Trình bày trung thực cả các điểm yếu có
thể ảnh hưởng đến kết luận.}

[Các hạn chế đã xác định và tác động có thể có.]

% ============================================================
\chapter{Kết luận và hướng phát triển}
% ============================================================

\section{Kết luận}

\huongdan{Trả lời trực tiếp từng câu hỏi nghiên cứu. Tóm tắt đóng góp chính,
kết quả quan trọng nhất và mức độ đạt được mục tiêu. Chỉ kết luận những gì
được dữ liệu và thí nghiệm hỗ trợ; không đưa kết quả mới chưa trình bày trước đó.}

[Những kết luận chính của nghiên cứu.]

\section{Đóng góp của đề tài}

\huongdan{Nêu rõ phần mới hoặc hữu ích do dự án thực hiện: phương pháp, quy trình,
hệ thống, bộ dữ liệu, phân tích hoặc bằng chứng thực nghiệm. Phân biệt đóng góp
của nhóm với các công cụ và kết quả kế thừa từ công trình trước.}

[Các đóng góp cụ thể và bằng chứng hỗ trợ.]

\section{Hướng phát triển}

\huongdan{Đề xuất những bước tiếp theo dựa trên hạn chế đã phát hiện hoặc câu hỏi
mới nảy sinh từ kết quả. Ưu tiên hướng phát triển cụ thể, có lý do, thay vì chỉ
nêu chung chung rằng sẽ ``cải thiện hệ thống''.}

[Các cải tiến hoặc thí nghiệm tiếp theo.]

% ============================================================
% TÀI LIỆU THAM KHẢO
% ============================================================
\begin{thebibliography}{99}

\bibitem{ref1}
[Tác giả].
\newblock [Tên bài báo hoặc tài liệu].
\newblock [Tên tạp chí/hội nghị/nhà xuất bản], [năm].
\newblock [DOI hoặc URL nếu có].

\bibitem{ref2}
[Tác giả].
\newblock [Tên bài báo hoặc tài liệu thứ hai].
\newblock [Nguồn xuất bản], [năm].
\newblock [DOI hoặc URL nếu có].

\end{thebibliography}

% ============================================================
% PHỤ LỤC (tùy chọn)
% Có thể đưa vào đây bảng hỏi, thông số chi tiết, giả mã, bảng số liệu
% đầy đủ, hướng dẫn tái lập thí nghiệm hoặc hình bổ sung.
% ============================================================
\appendix
\chapter{Thông tin bổ sung (nếu cần)}

\huongdan{Chỉ thêm phụ lục khi có tài liệu hỗ trợ việc kiểm chứng hoặc tái lập
nghiên cứu nhưng quá dài để đặt trong phần chính.}

[Thông tin bổ sung.]

\end{document}
